% Reference: https://en.cppreference.com/w/cpp/string/basic_string.html
\subsection{string：截取、查找与修改}
\label{stl:string}
\subsubsection*{原理}
\InlineCode{string} 管理连续字符序列，可用下标、迭代器及排序等算法。下标从 0 开始；字符串比较按字符字典序，而不是把数字串转换成数值比较。字符访问为 $O(1)$，截取与中间修改通常需要线性复制或移动。
\subsubsection*{使用教程}
\InlineCode{substr(pos,len)} 从 \InlineCode{pos} 开始复制至多 \InlineCode{len} 个字符；\InlineCode{find(t,pos)} 从指定位置查找首次出现，省略位置从 0 开始，未找到返回 \InlineCode{string::npos}。\InlineCode{erase(pos,len)} 删除，\InlineCode{insert(pos,t)} 插入，\InlineCode{+=} 追加，\InlineCode{clear()} 清空。
\begin{lstlisting}[style=mtycpp,numbers=none]
string s="abacaba";
string t=s.substr(1,3);
auto pos=s.find("aca");
if(pos!=string::npos) s.erase(pos,3);
s.insert(1,"xy");
s+='z';
reverse(s.begin(),s.end());
string num=to_string(123456789LL);
long long x=stoll(num);
\end{lstlisting}
\InlineCode{substr} 的第二项是长度，不是右端点；起点可以等于长度，但不能更大。不要先把 \InlineCode{npos} 转为 \InlineCode{int} 再判断。\InlineCode{stoll} 要求输入可解析且结果在 64 位范围内，否则可能抛异常；\InlineCode{cin>>s} 读取一个非空白单词，含空格的一行用 \InlineCode{getline}。
